\ifdefined\LFSWebBook\else
\documentclass[../textbook.tex]{subfiles}
\begin{document}
\fi
\bookchapter{自回归语言建模}{ch:rev-autoregressive}

一个语言模型需要同时回答两个不同问题：如何为序列分配概率，以及如何根据这些概率产生实际输出。前者由模型的条件分布与训练目标定义，后者由解码策略与终止规则定义。相同权重可以在贪心选择、随机采样和束搜索下产生不同文本；较低的预测损失与自由生成时的重复、遗漏或事实错误也各有来源。

自回归生成将条件概率、解码选择与终止规则连接起来。因果网络计算下一词元分布，教师强制使用已知目标前缀训练网络，解码器则在生成时逐步选择词元。EOS、停止串和长度预算共同规定输出的终止方式。

\section{用条件分布定义序列}

\subsection{自回归分解的结构范围}
设词表为有限集合 $\mathcal V$，词表大小为 $V_{\mathrm{vocab}}$。对长度为 $T$ 的词元序列 $x_{1:T}$，概率链式法则给出
\begin{equation}
 p_\theta(x_{1:T})=\prod_{t=1}^{T}p_\theta(x_t\mid x_{<t}),\qquad
 \log p_\theta(x_{1:T})=\sum_{t=1}^{T}\log p_\theta(x_t\mid x_{<t}).\label{eq:ar-chain}
\end{equation}
\term{自回归}{Autoregression}{}表示后续变量的分布以已经给定的变量为条件。链式分解本身不要求使用神经网络；限定上下文长度、共享各位置参数、采用因果 Transformer 才是具体建模选择。

对于给定提示或源输入 $c$ 的条件生成，同样有
\begin{equation}
 p_\theta(y_{1:T}\mid c)=\prod_{t=1}^{T}p_\theta(y_t\mid c,y_{<t}).\label{eq:ar-conditional}
\end{equation}
条件 $c$ 可以是同一序列中的提示前缀，也可以通过编码器及交叉注意力提供。两种方式都能实现条件生成，条件概率分解与具体网络模块相互独立。

有限上下文模型实际计算的可能是 $p_\theta(x_t\mid x_{\max(1,t-C):t-1})$，其中 $C$ 是可读取的上下文长度。它对完整历史施加了信息限制，仍可逐步构成一致的序列分布，但不再保证利用更早内容。裁掉前缀或改变位置编号会改变模型的条件，其影响超出内存占用。

\begin{table}[htbp]
\centering
\caption{自回归建模与解码生成核心数学符号}\label{tab:rev-generation-notations}
\begin{tabularx}{\linewidth}{p{24mm}X}
\toprule
符号 & 含义与约束\\\midrule
$x_t,y_t,c$ & 观测词元、生成词元与外部条件；词元取值属于 $\mathcal V$。\\
$B,T,C$ & 批量大小、序列位置数、上下文容量；不使用 $T$ 表示温度。\\
$d,V_{\mathrm{vocab}}$ & 隐藏宽度与词表大小。\\
$z_t,p_t,q_t$ & 当前词表分数、原始模型分布、经过解码变换的实际分布。\\
$\tau,k,\rho$ & 正温度、Top-$k$ 个数、Top-$p$ 累积阈值；$\tau>0$，$1\leq k\leq V_{\mathrm{vocab}}$，$0<\rho\leq1$。\\
$K,L_{\max}$ & 束宽与最多新增词元数，均为正整数。\\
\bottomrule
\end{tabularx}
\end{table}

\subsection{可变长度序列与结束事件}
仅对固定长度写式\eqref{eq:ar-chain}，尚未说明文本在哪里结束。引入\term{序列结束符}{End of Sequence}{EOS}，将一条完整响应定义为非结束词元 $y_{1:T}$ 后跟 EOS：
\begin{equation}
 p_\theta(y_{1:T},\mathrm{EOS}\mid c)
 =\left[\prod_{t=1}^{T}p_\theta(y_t\mid c,y_{<t})\right]
 p_\theta(\mathrm{EOS}\mid c,y_{1:T}).\label{eq:ar-eos-prob}
\end{equation}
结束概率也是序列概率的一部分。比较两个完整候选时，漏掉 EOS 会改变排名。初始条件可以由提示前缀给出，也可以借助\term{序列开始符}{Beginning of Sequence}{BOS}建立；不少模型并不使用单独的 BOS 词元。

各步分布归一化只约束当前步候选的概率和。有限时间终止还取决于 EOS 概率；例如某类前缀之后 EOS 概率恒为零，就可能产生永不终止的路径。如果每个未结束前缀上的 EOS 概率都有统一下界 $\epsilon>0$，那么
\begin{equation}
 \Pr(\text{前 }t\text{ 步仍未结束})\leq(1-\epsilon)^t\longrightarrow0,
\end{equation}
这是确保几乎必然终止的一种充分条件，真实模型与截断解码往往不满足。实际系统因此仍需要明确的有限生成预算。

\section{下一词元分布}

\subsection{GPT 数据流}
\term{生成式预训练 Transformer}{Generative Pre-trained Transformer}{GPT}以仅解码器的因果网络为常见结构。输入编号先经嵌入与位置处理变成隐藏向量，再通过多层因果注意力、逐位置前馈变换、残差及归一化，最后映射为词表上的实数分数。具体模型的归一化、位置方案与前馈形式属于结构配置，GPT 这一缩写本身不足以确定这些细节。

取输入前缀 $x_{1:t}$，末位隐藏行向量为 $\vect{h}_t\in\Real^{1\times d}$。词表输出头计算
\begin{equation}
 \vect{z}_t=\vect{h}_t\mat{W}_{\mathrm{lm}}+\vect{b},\qquad
 \vect{p}_\theta(x_{t+1}=v\mid x_{\leq t})
 =\frac{\exp z_{t,v}}{\sum_{u\in\mathcal V}\exp \vect{z}_{t,u}},\label{eq:ar-head}
\end{equation}
其中 $\mat{W}_{\mathrm{lm}}\in\Real^{d\times V_{\mathrm{vocab}}}$，$\vect{b}\in\Real^{V_{\mathrm{vocab}}}$ 可按结构省略。$\vect{z}_t$ 是\term{未归一化分数}{Logits}{}，与概率相差一次 Softmax。若输入嵌入表为 $\mat{E}\in\Real^{V_{\mathrm{vocab}}\times d}$，权重共享可设置 $\mat{W}_{\mathrm{lm}}=\mat{E}^{\mathrm{T}}$，但这是显式参数约束，维度相容本身推不出它。

\begin{figure}[htbp]
\centering
\begin{tikzpicture}[>=Stealth,node distance=7mm,
 block/.style={draw,rounded corners,align=center,font=\small,text width=29mm,minimum height=11mm}]
 \node[block] (id) {词元前缀\\$[B,T]$};
 \node[block,right=of id] (emb) {嵌入与位置\\$[B,T,d]$};
 \node[block,right=of emb] (net) {因果网络堆叠\\$[B,T,d]$};
 \node[block,right=of net] (head) {归一化与输出头\\$[B,T,V_{\mathrm{vocab}}]$};
 \draw[->] (id)--(emb);\draw[->] (emb)--(net);\draw[->] (net)--(head);
 \node[block,below=of head] (loss) {训练：各有效位置\\匹配下一个真实词元};
 \node[block,below=of net] (next) {生成：末位置分布\\选择一个新词元};
 \draw[->] (head)--(loss);\draw[->] (head.south west)--(next.north east);
\end{tikzpicture}
\caption{相同因果网络的两种使用方式。训练利用所有有效位置的已知目标，生成读取当前前缀末位的分布。}
\label{fig:ar-gpt-flow}
\end{figure}

\subsection{因果性与前缀一致性}
\bookxref{ch:attention}[定义的因果注意力]通过因果掩码使位置 $t$ 只能读取不晚于 $t$ 的信息。如果每层其他操作均逐位置执行，则完整序列前向在位置 $t$ 得到的条件分数，与只输入前缀 $x_{1:t}$ 时末位的分数在数学上相同。

这个结论需要相同的参数、位置编号、有效性掩码和运行模式。随机失活开启时，两次调用可能使用不同随机掩码；浮点内核也可能因张量形状不同产生舍入差异。上述因素与理论上的未来信息泄漏属于不同层次。因果掩码不会自动修正标签对齐错误，预测对象还需要单独定义。

\section{教师强制与移位监督}

\subsection{输入目标错位}
\term{教师强制}{Teacher Forcing}{}在训练时始终给定数据中的真实前缀，再预测其后一个真实词元。设原序列为
$(\mathrm{BOS},\text{我},\text{学习},\text{模型},\mathrm{EOS})$，一个明确的训练对为
\begin{equation}
 \begin{aligned}
 \text{输入}&=(\mathrm{BOS},\text{我},\text{学习},\text{模型}),\\
 \text{目标}&=(\text{我},\text{学习},\text{模型},\mathrm{EOS}).
 \end{aligned}
\end{equation}
每个输入位置可以读取自身及过去输入，但其目标是下一位置的词元。为避免“左移”“右移”因参照物不同引起歧义，应直接说明索引关系：若输入位置 $t$ 存放 $x_t$，则该位置分数监督 $x_{t+1}$。构造解码器输入时也可以说“在目标前面放入 BOS，再移除最后一个词元”，二者是同一个对齐关系。

\begin{figure}[htbp]
\centering
\begin{tikzpicture}[>=Stealth,x=2.25cm,y=1cm]
 \foreach \x/\a/\b in {0/BOS/我,1/我/学习,2/学习/模型,3/模型/EOS}{
  \node[draw,minimum width=17mm,minimum height=8mm] (in\x) at (\x,0) {\a};
  \node[draw,minimum width=17mm,minimum height=8mm] (out\x) at (\x,-1.6) {\b};
  \draw[->] (in\x)--node[right,font=\small] {预测}(out\x);
 }
 \node[left,font=\small] at (-.5,0) {输入};
 \node[left,font=\small] at (-.5,-1.6) {目标};
 \draw[->,dashed] (in0.east)--(in1.west);
 \draw[->,dashed] (in1.east)--(in2.west);
 \draw[->,dashed] (in2.east)--(in3.west);
\end{tikzpicture}
\caption{监督索引与因果前缀。竖直箭头给出各位置的预测对象；横向虚线表示前缀依赖，实际注意力仍读取全部允许历史。}
\label{fig:ar-shift}
\end{figure}

移位可以由数据整理器完成，也可以由模型损失内部完成。前一种接口接收已经对齐的输入和目标；后一种接口可能接收相同的完整词元数组，再在内部使用前 $T-1$ 个分数与后 $T-1$ 个目标。两处同时移位会变成跨两步预测，两处都不移位则可能学到复制当前位置。应以明确索引式作为协议，参数名无法替代它。

\subsection{平均负对数似然}
给定训练文档 $x^{(b)}$，最大化其序列似然等价于最小化各条件负对数概率之和。设 $m_{bt}\in\{0,1\}$ 标记位置 $(b,t)$ 是否提供有效监督，则
\begin{equation}
 \mathcal L(\theta)=-\frac1N\sum_{b,t}m_{bt}
 \log p_\theta(x^{(b)}_{t+1}\mid x^{(b)}_{\leq t}),\qquad
 N=\sum_{b,t}m_{bt}>0.\label{eq:ar-loss}
\end{equation}
填充目标通常不计入 $N$，文档末尾的 EOS 若属于训练目标则计入。指令响应训练中，提示位置可以参与条件计算而不承担损失；这相当于选择监督区域，屏蔽所有提示键的做法与之不同。

若批次含不同有效长度，应聚合负对数似然总和后除以有效目标总数。两个批次有效目标数分别为 $10$ 与 $90$、平均损失分别为 $1$ 与 $3$ 时，整体平均是 $\frac{10\cdot1+90\cdot3}{100}=2.8$，与逐批平均后再取平均得到的 $\frac{1+3}{2}=2$ 不同。按序列平均和按词元平均也是不同的目标权重设计。

在使用自然对数时，\term{困惑度}{Perplexity}{PPL}为 $\exp\mathcal L$。它概括同一数据和分词规则下的条件预测损失，与模型候选答案数量无关，也无法跨不同分词器直接比较。

\subsection{因果训练并行性}
训练样本已经给出全部真实目标，所以位置 $1,\ldots,T$ 所需的合法前缀同时可用。因果掩码阻止未来泄漏，矩阵运算仍可在同一次网络调用中处理全部位置。各网络层之间保留先后依赖，这里的并行指位置维度的批量计算，与所有运算在时间上同时完成不同。

自由生成时，模型还不知道 $y_t$，而计算 $p(y_{t+1}\mid c,y_{\leq t})$ 需要先确定它。生成因此沿实际选出的前缀逐步推进。推测解码可以批量提出并验证候选，成为后续条件的只有保留的前缀，自回归分解本身依然成立。

\begin{algorithm}[移位目标的教师强制损失]
\label{alg:ar-teacher}
\textbf{输入：}各文档的词元序列、文档边界、填充标记和监督区域。\textbf{输出：}标量损失与有效目标数。\textbf{状态：}当前固定参数 $\theta$；本算法只计算一次目标，不推进优化器。
\begin{lfsalgorithmic}[1]
 \State 在每个合法文档范围内建立输入位置 $x_t$ 与目标 $x_{t+1}$ 的对应关系；依据是否包含 BOS、EOS 明确首尾监督。
 \State 构造因果和填充可见性，标记有效目标 $m_t$。跨文档是否连接必须遵守事先定义的目标，不在批次边界临时决定。
 \State 运行因果网络，得到与输入位置对应的分数；确认移位仅在数据侧或损失侧执行一次。
 \State 累加有效位置的负对数概率和以及目标数 $N$；$N=0$ 时返回明确的无监督批次状态，不执行除零或伪造零梯度。
 \State 返回总和除以 $N$。各目标只能读取其前缀，填充位置不贡献监督，输入和目标索引保持图\ref{fig:ar-shift}的关系。
\end{lfsalgorithmic}
\end{algorithm}

\section{真实前缀与生成前缀}

\subsection{暴露偏差的含义与限度}
训练时前缀来自数据分布，生成时前缀来自模型和解码策略。有限模型在少见或偏离训练分布的前缀上可能预测较差，一步错误又会成为下一步的条件。这种训练与使用条件的差别通常称为\term{暴露偏差}{Exposure Bias}{}，是序列训练研究的重要动机之一\cite{bengio2015scheduledsampling}。

教师强制直接计算序列的负对数似然。对真实序列分布 $P$ 和模型分布 $p_\theta$，在支持集及可积性条件满足时，
\begin{equation}
 \mathbb E_{x\sim P}[-\log p_\theta(x)]
 =H(P)+D_{\mathrm{KL}}(P\|p_\theta).
\end{equation}
因序列链式法则，教师强制恰好计算这一序列负对数似然。若模型准确表达了真实条件分布，其生成分布也可以一致。实际困难来自有限数据、模型失配、分布外前缀，以及最终任务指标与局部概率目标之间的差异。
人为混入模型预测前缀，会改变训练数据条件，并且仍需判断原目标是否与被修改的前缀相容。修改后的前缀与目标需要保持语义一致。序列级训练、偏好优化和可验证奖励从不同角度调整行为，后训练篇将讨论其目标与条件。

\subsection{生成误差累积}
如果在一组明确的条件下，每一步出现某类错误的条件概率不超过 $\epsilon$，则由并集界，前 $T$ 步至少一次错误的概率不超过 $T\epsilon$，且上界至多取一。在独立且等概率的特殊模型中，该概率为 $1-(1-\epsilon)^T$。真实生成中的错误相互影响，独立同分布假设通常不成立；逐步错误上界需要按生成前缀测量，教师强制平均损失给出的是另一种总体统计量。

需要区分预测某个指定词元的概率与完成任务的正确率。某个问题可以有许多正确表述；合法同义词的概率较低未必是任务错误，重复的高概率短句也未必完成了任务。自由生成评价应同时关注内容、约束、终止及整体连贯性。

\section{幻觉、事实性与证据支持}

\subsection{事实性判定基准}
\term{幻觉}{Hallucination}{}通常指生成结果包含相对于任务规定参照不受支持或错误的内容\cite{ji2023hallucination}。这个词描述输出与参照之间的可检验关系，与模型是否具有人的感知、信念或欺骗意图无关。参照可以是指定上下文、带版本与时间的外部事实、工具实际返回值，或任务明确给出的约束；若不先说明参照，“是否幻觉”便可能没有唯一判定。

例如，某个结论可能碰巧符合外部世界，却没有得到当前证据包支持；它在答案正确性上可能通过，在证据忠实度上则失败。反过来，模型也可能忠实复述一份已经过期的材料，因而有证据支持却不符合问题要求的当前事实。这两个维度各有判据，合并为一个模糊的“看起来正确”会丢失区分。

\begin{table}[htbp]
\centering
\caption{大语言模型生成事实性异常分类与判定基准。}
\label{tab:ar-factuality-criteria}
\begin{tabularx}{\linewidth}{@{}lX@{}}
\toprule 现象 & 相对于已声明参照的判定\\\midrule
事实错误 & 可核实声明与目标时间、适用范围内的有效事实冲突。\\
上下文矛盾 & 回答声明与给定证据直接冲突，或忽略证据中的必要限定。\\
无依据生成 & 回答新增了可核实声明，但给定证据既不支持它，也不足以推出它。\\
虚构归因 & 回答捏造来源、引用、工具结果或把真实来源错误地说成支持某一声明。\\
\bottomrule
\end{tabularx}
\end{table}

这些类型可以重叠：虚构引用所支持的结论也可能同时是事实错误。漏答、文风不佳、格式错误和正确拒答属于各自的任务维度，一个宽泛的“幻觉”术语不应吞并所有生成失败。

\subsection{词元概率与事实性}
模型训练直接约束的是式\eqref{eq:ar-loss}中的下一词元条件概率，每个自然语言声明的真值不在这一目标之内。对回答中的声明 $A_i$、证据集合 $E$ 和目标时间 $t^*$，模型给出的
\begin{equation}
 p_\theta(y_j\mid c,y_{<j})
 \quad\text{与}\quad
 \Pr(A_i\text{ 在 }E,t^*\text{ 下成立})
\end{equation}
是不同对象，前者的数值不构成后者的校准估计。一个常见而流畅的错误说法可以由高概率词元组成；一个罕见但正确的名称也可能具有较低词元概率。

幻觉可能来自训练语料的遗漏、冲突或过时，有限模型与优化误差，问题歧义或证据缺失，以及生成前缀偏离后形成的误差反馈。后训练若主要奖励完整、肯定和流畅的回答，也可能使模型在证据不足时仍倾向作答。外部检索又会引入漏召回、错误排序、版本冲突和证据未被正确利用等系统原因。观察到错误答案后，仅凭文本无法断定故障发生在预训练、解码器还是检索层。

降低温度会集中已有分布，但不会增加知识或执行事实核验；若最高概率路径本来就是错误的，贪心或低温只会更稳定地返回它。提高温度则可能增加输出变化和低概率错误，后者与幻觉也各自独立。事实核验需要可用证据或外部验证。

\subsection{幻觉缓解机制}
减少幻觉应把生成系统拆成可检查阶段。首先规定允许使用的知识范围、目标时间和不可回答条件；其次检索带来源、版本与权限的证据，并在生成时限制结论范围、保留声明到来源的映射；随后对关键声明执行语义支持检查、规则检查、计算或工具验证；证据不足或冲突无法解决时，输出限定回答、部分回答或拒答。RAG 可以提供新鲜证据，但漏召回、错误证据或生成器忽略证据时，幻觉依然可能存在，完整机制见\bookxref{ch:rag-systems}。

评估时应同时包含可回答、不可回答和证据冲突样例，把回答拆成可核实声明，并分别标记支持、冲突与未知。答案正确率、声明级忠实度、引用精确率、正确拒答率和过度拒答率的分母不同，应分别报告；具体测量方法见\bookxref{ch:evaluation-statistics}。平均语言模型损失和词面相似度衡量各自定义的预测与文本差异，幻觉率需要声明级事实与证据标注。

\begin{note}[流畅生成与可靠断言是两项能力]
自回归模型可以对语言形式给出连贯的条件分布，事实判断还需要明确参照、可用证据、时间范围和验证规则。降低随机性可以减少答案波动；事实可靠性由证据与验证规则确定。系统应在证据不足或冲突时输出限定结论或拒答。\end{note}

\section{采样分布变换}

\subsection{温度的概率形式}
\term{温度缩放}{Temperature Scaling}{}在生成时使用 $\tau>0$ 改变分数尺度：
\begin{equation}
 p^{(\tau)}_v=\frac{\exp(\frac{z_v}{\tau})}{\sum_u\exp(\frac{z_u}{\tau})}
 =\frac{p_v^{\frac{1}{\tau}}}{\sum_up_u^{\frac{1}{\tau}}}.\label{eq:ar-temperature}
\end{equation}
温度调整生成分布的集中程度。正温度不改变分数排序，但改变候选间相对概率：
\begin{equation}
 \frac{p^{(\tau)}_i}{p^{(\tau)}_j}
 =\exp\!\left(\frac{z_i-z_j}{\tau}\right).
\end{equation}
当 $\tau\to0^+$ 且最大分数唯一时，分布趋于该候选上的点质量；若多个有限最大值相同，则趋于这些最大值上的均匀分布。$\tau=0$ 会使除法失去定义，若要贪心选择，应采用明确的离散分支。$\tau\to\infty$ 时，在有限分数支持集上趋于均匀。已被硬屏蔽的候选不因此恢复。

\subsection{温度与熵的关系}
\begin{assumption}[温度单调性的适用条件]
本节固定同一个前缀、同一组有限分数和非空候选支持集，只改变正温度。求导期间不重新训练、不修改截断集合，也不切换到另一个采样前缀。概率与熵均指这一固定条件下的分布。
\end{assumption}
对固定有限分数集合，令 $\beta=\frac{1}{\tau}$、$Z(\beta)=\sum_v \conste{}^{\beta z_v}$。熵写为
\begin{equation}
 H(\beta)=\log Z(\beta)-\beta\mathbb E_{p^{(\tau)}}[\vect{z}].
\end{equation}
由 $\frac{\dif\log Z}{\dif\beta}=\mathbb E[\vect{z}]$ 以及
$\frac{\dif\mathbb E[\vect{z}]}{\dif\beta}=\operatorname{Var}(\vect{z})$，有
\begin{equation}
 \frac{\dif H}{\dif\beta}=-\beta\operatorname{Var}(\vect{z}),\qquad
 \frac{\dif H}{\dif\tau}=\frac{\operatorname{Var}_{p^{(\tau)}}(\vect{z})}{\tau^3}\geq0.
\end{equation}
因此在固定支持集与固定分数下，提高温度不会降低条件熵。这是关于同一时刻分布的结果；不同采样路径会产生不同前缀，最终文本的质量与多样性需在完整生成过程上测量。

\subsection{Top-k 与核采样}
\term{前 $k$ 项采样}{Top-k Sampling}{}保留概率最大的 $k$ 个候选，令集合为 $\mathcal K$，再归一化：
\begin{equation}
 q_v=\frac{p_v\mathbf1[v\in\mathcal K]}{\sum_{u\in\mathcal K}p_u}.
\end{equation}
边界并列时可以按稳定规则恰好选择 $k$ 项，也可以保留所有达到阈值的项；后一种方式的支持集可能多于 $k$。两种定义均应明确，同一个配置名之下的实际集合可能不同。

\term{核采样}{Nucleus Sampling}{Top-p}按概率从大到小排序，取累计概率至少达到 $\rho$ 的最小前缀集合\cite{holtzman2020degeneration}。若排序概率为 $p_{(1)}\geq\cdots\geq p_{(V_{\mathrm{vocab}})}$，则
\begin{equation}
 m=\min\left\{r:\sum_{j=1}^{r}p_{(j)}\geq\rho\right\},\qquad
 \mathcal N_\rho=\{(1),\ldots,(m)\}.
\end{equation}
之后在 $\mathcal N_\rho$ 内重新归一化并采样。跨过阈值的那个候选必须保留，尤其在第一候选就超过阈值时，集合仍至少包含一个词元。Top-p 中的 $p$ 是方法名称，正文用 $\rho$ 表示阈值，以区别模型概率。

Top-k 固定候选数量，Top-p 固定保留的累计概率门槛。尖锐分布下 Top-p 可能只留下一个词元，平坦分布下则可能留下很多。二者截断概率质量的几何差异如图\ref{fig:sampling-geometry-truncation}所示。二者均改变支持集，可能移除 EOS 或任务所需的低概率候选，质量上的代价需要在具体任务中衡量。

\begin{figure}[htbp]
\centering
\includegraphics[width=0.96\linewidth]{\subfix{../../figures/theory/sampling-geometry-truncation.pdf}}
\caption{Top-$k$ 截断与 Top-$p$ 核采样的几何选择机制对比。左图展示固定候选数 $k$ 的硬截断；右图展示依据动态累积概率质量 $\rho$ 的自适应核集合划分。}\label{fig:sampling-geometry-truncation}
\end{figure}

\subsection{采样变换顺序}
\begin{example}
设某个前缀后的原始概率为
\begin{equation}
 p=(0.4,0.3,0.2,0.1),
\end{equation}
对应分数可取 $z=(\log0.4,\log0.3,\log0.2,\log0.1)$。温度 $\tau=\frac{1}{2}$ 相当于将概率平方再归一化，平方和为 $0.30$，所以
\begin{equation}
 p^{(\frac{1}{2})}=\left(\frac8{15},\frac3{10},\frac2{15},\frac1{30}\right).
\end{equation}
原分布做 Top-$2$ 得到 $(\frac{4}{7},\frac{3}{7},0,0)$；原分布做阈值 $\rho=0.75$ 的 Top-p 时，前两项之和为 $0.7$，必须保留第三项，得到 $(\frac{4}{9},\frac{1}{3},\frac{2}{9},0)$。

若先降温再做同样的 Top-p，前两项的质量为 $\frac{8}{15}+\frac{3}{10}=\frac{5}{6}>0.75$，只保留两项，最终分布为 $(\frac{16}{25},\frac{9}{25},0,0)$。如果先在原分布上做 Top-p，再降温，支持集仍有三项，最终为 $(\frac{16}{29},\frac{9}{29},\frac{4}{29},0)$。因此温度与 Top-p 的次序会改变分布。温度不改变候选排序，所以在没有其他变换、并列规则相同的条件下，它与固定 Top-k 支持集选择的关系不同；一个特例不足以推出所有处理器都可交换。
\end{example}

\subsection{生成分布偏移}
设经过温度、截断和其他规则后的条件分布为 $q_t(\cdot\mid c,y_{<t})$，实际随机生成的路径概率为
\begin{equation}
 q(y_{1:T}\mid c)=\prod_{t=1}^{T}q_t(y_t\mid c,y_{<t}).
\end{equation}
它通常不等于原始 $p_\theta(y_{1:T}\mid c)$。即使只用相同温度，每个前缀的归一化常数不同，把结果解释为“整条序列概率统一取 $\frac{1}{\tau}$ 次方后归一化”一般也不成立。逐步变换定义的是另一组条件分布。

记录生成概率、计算重要性比率或检验精确采样时，必须说明使用原始模型分布还是实际解码分布。温度和截断不会改变模型权重，也不会为模型补充事实知识；它们改变的是已给定分数被转化为决策的方式。

\section{贪心与束搜索}

\subsection{局部最优与序列最优}
\begin{example}
\term{贪心解码}{Greedy Decoding}{}每一步选择条件概率最大的候选，按固定规则解决并列。它最大化当前一步，序列层面的总分与之无关。对完整序列，应比较式\eqref{eq:ar-eos-prob}中的乘积，计算时通常累加对数。

构造一个有限分支概率模型：第一步 $p(a)=0.6,p(b)=0.4$；在 $a$ 后，EOS 与 $c$ 的概率各为 $0.5$；在 $b$ 后，EOS 概率为 $0.9$、$c$ 为 $0.1$；一旦产生 $c$，下一步必为 EOS。其完整序列分布为
\begin{table}[htbp]
\centering
\caption{贪心解码局部最优与序列全局最优不一致示例。}\label{tab:rev-greedy-failure}
\begin{tabular}{lc}
\toprule
完整序列 & 概率\\\midrule
$a,\mathrm{EOS}$ & $0.6\cdot0.5=0.30$\\
$a,c,\mathrm{EOS}$ & $0.6\cdot0.5\cdot1=0.30$\\
$b,\mathrm{EOS}$ & $0.4\cdot0.9=0.36$\\
$b,c,\mathrm{EOS}$ & $0.4\cdot0.1\cdot1=0.04$\\
\bottomrule
\end{tabular}
\end{table}
四条路径概率之和为一。贪心第一步选 $a$，无论随后并列选择 EOS 还是 $c$，完整概率均为 $0.30$，低于 $b,\mathrm{EOS}$ 的 $0.36$。未列出的分支概率设为零。

\begin{figure}[htbp]
\centering
\begin{tikzpicture}[>=Stealth,node distance=12mm,
 nd/.style={draw,rounded corners,font=\small,align=center,minimum width=19mm,minimum height=8mm}]
 \node[nd] (root) {起点};
 \node[nd,right=20mm of root,yshift=13mm] (a) {$a:0.6$};
 \node[nd,right=20mm of root,yshift=-13mm] (b) {$b:0.4$};
 \node[nd,right=26mm of a,yshift=7mm] (ae) {$a,\mathrm{EOS}:0.30$};
 \node[nd,right=26mm of a,yshift=-7mm] (ac) {$a,c:0.30$};
 \node[nd,right=26mm of b,yshift=7mm] (be) {$b,\mathrm{EOS}:0.36$};
 \node[nd,right=26mm of b,yshift=-7mm] (bc) {$b,c:0.04$};
 \draw[->] (root)--(a);\draw[->] (root)--(b);
 \draw[->] (a)--(ae);\draw[->] (a)--(ac);
 \draw[->] (b)--(be);\draw[->] (b)--(bc);
\end{tikzpicture}
\caption{局部贪心错过完整高概率序列的概率树。节点数值为累计概率；含 $c$ 的前缀下一步以概率一结束。}
\label{fig:ar-beam-tree}
\end{figure}
\end{example}

\subsection{束搜索保留多个前缀}
\term{束搜索}{Beam Search}{}以有限宽度 $K$ 保留候选，反复扩展与剪枝。设前缀分数为
\begin{equation}
 s(y_{1:t})=\sum_{j=1}^{t}\log p_\theta(y_j\mid c,y_{<j}).\label{eq:ar-beam-score}
\end{equation}
每次扩展一个候选 $v$，只需令 $s' = s+\log p_\theta(v\mid c,y_{1:t})$。同一深度的候选可以批量前向，但各自前缀、缓存和结束状态必须保持对应，排序后错配状态会导致错误结果。

图\ref{fig:ar-beam-tree}中宽度为二时，第一步保留 $a,b$；第二步最高的扩展是 $b,\mathrm{EOS}$，得分对应 $0.36$，其次为两条 $0.30$ 分支之一。因此保留多个前缀可修复该例的贪心选择。有限宽度仍可能过早删除一条未来更优的路径；增大束宽不提供一般的全局最优保证，任务质量也不会随之提升。

\begin{algorithm}[以原始对数概率评分的束搜索]
\label{alg:ar-beam}
\textbf{输入：}条件 $c$、束宽 $K$、EOS 集合、最大新增长度 $L_{\max}$ 和确定的并列规则。\textbf{输出：}一个完成候选及终止原因，或未完成候选及长度截断标记。\textbf{状态：}至多 $K$ 条记录，每条含前缀、累计分数、完成标记及其模型状态。
\begin{enumerate}
 \item 初始化空响应候选，累计分数为零，状态与条件 $c$ 对应。
 \item 对每条未完成候选计算下一步分布，为每个有限分数词元生成扩展并累计式\eqref{eq:ar-beam-score}；若词元属于 EOS 集合，将该扩展标记完成。已完成候选原样进入候选池，不再追加词元或重复计分。
 \item 从候选池按原始累计分数保留最高的 $K$ 条，连同其对应模型状态一起重排。若不足 $K$ 条则保留全部。
 \item 若当前最佳完成候选的分数不低于所有未完成前缀分数，可以停止；若全部候选完成，也停止。否则继续扩展，最多进行 $L_{\max}$ 次。
 \item 有完成候选时返回其中最高分者；没有时返回最高分未完成前缀并报告截断。外部预算耗尽与模型选择 EOS 分别记录为截断和正常终止。
\end{enumerate}
本算法让完成候选占据束内位置，定义了一种明确的有限束策略；$K=1$ 时与采用相同并列及终止规则的贪心解码一致。维护独立完成池的实现可以使用不同的候选预算，其规则需要单独说明。
\end{algorithm}

\subsection{束搜索停止条件}
每个条件概率不超过一，所以 $\log p\leq0$，任何未完成前缀的未来扩展分数都不超过当前分数。如果一个完成候选已经不低于所有保留前缀分数，那么继续扩展这些前缀不可能超过它。这证明了算法\ref{alg:ar-beam}中的停止条件，但仅相对于尚未剪去的候选成立，此前被删除的路径无法重新发现。

长度归一化会改变这个性质。常见目标可写为 $\frac{s(y)}{|y|^\alpha}$，$\alpha>0$，其中长度是否包含 EOS 必须固定。因为 $s(y)$ 为负，分母随长度增大可能使归一化分数上升，当前前缀分数不再是未来得分的上界。若存在确定的最大长度，可以根据该目标另行构造上界；若还加入奖励或覆盖项，则必须按完整评分重新推导，原始对数概率的停止证明不再适用。

\subsection{长度偏好与任务目标}
较长路径累加更多非正项，原始序列概率常偏向某些短输出，但这不等于“序列越长概率一定更低”的跨路径结论，因为不同路径的条件概率不同。长度归一化、最小长度与重复惩罚可以改变搜索行为，同时也改变目标，它们服务的任务约束需要写明。

最大概率输出与高质量输出不是同一个数学目标。开放式任务有大量合理续写，高概率路径可能重复或过于笼统；有严格约束的翻译与结构化任务也可能因长度或覆盖问题失真。解码选择应围绕任务评价，束宽增大与模型能力提升是两回事。核采样相关研究讨论的退化现象也有特定模型和任务条件，无法外推成所有任务上唯一最优策略\cite{holtzman2020degeneration}。表\ref{tab:rev-decoding-comparison}全面对比了自回归生成中的主流解码与采样策略。

\begin{table}[htbp]
\centering
\caption{自回归语言模型主流解码与采样策略全景对比}\label{tab:rev-decoding-comparison}
\begin{tabularx}{\linewidth}{p{18mm}p{22mm}p{18mm}p{22mm}X}
\toprule
解码策略 & 核心决策机制 & 随机性与多样性 & 计算复杂度与延迟 & 典型适用场景与退化风险 \\\midrule
贪心解码 & $\arg\max_v p(v)$ 每步取极值 & 确定性输出（零随机性） & $\mathcal{O}(L)$ 单步单前缀，延迟最低 & 事实问答、代码补全、数学推理；长文本生成易陷入循环复读 \\
束搜索 (Beam) & 维护 $K$ 条最高累计对数概率路径 & 弱随机性（可设束宽惩罚） & $\mathcal{O}(K \cdot L)$ 每步扩展 $K$ 倍前缀 & 机器翻译、语音识别、文本摘要；开放对话易偏向平淡通用短句 \\
温度采样 & $p_v^{(\tau)} \propto \exp(z_v/\tau)$ 重标度 & 可调；$\tau\to0$ 退化为贪心 & $\mathcal{O}(L)$，额外除法与 Softmax & 通用文本创作；高温易产生幻觉无意义词，低温缺乏多样性 \\
Top-$k$ 采样 & 硬截断前 $k$ 个最大概率词元 & 限制在 $k$ 个候选内随机采样 & $\mathcal{O}(L + V\log k)$ Top-$k$ 堆排序 & 故事续写、对话生成；分布平坦时截断过重，尖锐时引入噪声 \\
Top-$p$ (核采样) & 动态截取累积质量达 $\rho$ 的最小集合 & 动态自适应调节候选集大小 & $\mathcal{O}(L + V\log V)$ 全词表排序 & 现代主流采样方案；有效过滤长尾低质词同时保留多样性 \\
Min-$p$ 采样 & 截取概率大于 $p_{\max}\cdot p_{\text{base}}$ 候选 & 动态基于顶端置信度过滤 & $\mathcal{O}(L)$ 无需全排序，单遍扫描 & 平衡严谨逻辑与发散创造，防止高置信场景混入低质无关词 \\
推测解码 & 小模型草稿预测 $+$ 大模型并行验证 & 等价于目标大模型采样分布 & 延迟降低 $2\sim3\times$，加速自回归推理 & 大规模在线服务推理加速；要求草稿模型接受率较高 \\
\bottomrule
\end{tabularx}
\end{table}

\section{终止、约束与重复退化}

\subsection{生成停止条件}
模型采样出 EOS 表示模型分布选择结束；命中应用规定的停止串表示外部协议结束；耗尽 $L_{\max}$ 表示预算截断。三者的语义不同，应保留在返回状态中。EOS 通常不作为自然语言正文展示，但它仍属于模型的输出事件。

停止串可能跨词元边界，也可能跨流式分片边界；应对连续解码状态判断匹配，逐分片做独立字符串检查会漏掉匹配。是否保留停止串、匹配多个停止串时如何处理、特殊词元是否参与文本解码，都属于输出协议。\footnote{某些对话格式区分“一个消息结束”和“本轮助手回答结束”。这些标记即使都表现为特殊词元，也需要模型模板的说明才能确定能否合并为同一 EOS 语义。}

最小生成长度通常通过暂时禁止 EOS 来实现，它在允许长度之前改变分布。若用户请求短答案却强制较长最小长度，模型可能被迫继续补写。最大新增词元数与“提示加响应的最大总长度”也属于两项配置，前者控制输出预算，后者受到上下文容量约束，记录时必须区别。

\subsection{带明确状态的逐步生成}
\begin{algorithm}[温度与核采样的有限生成过程]
\label{alg:ar-sampling}
\textbf{输入：}已经序列化的条件词元、词表与模板、温度 $\tau>0$、核阈值 $\rho$、EOS 和停止串、最大新增长度 $L_{\max}$。\textbf{输出：}响应词元、文本与停止原因。\textbf{状态：}当前前缀、随机数状态、可选键值缓存以及连续文本解码状态。
\begin{lfsalgorithmic}[1]
\State 检查 $\tau>0$、$0<\rho\le1$、非负整数 $L_{\max}$ 和条件位置范围；固定推理模式
\State 当前前缀设为条件，响应置空；停止原因为长度截断
\For{$t=1,\ldots,L_{\max}$}
  \State 计算当前末位置分数，施加已声明的硬约束
  \If{允许候选集合为空}
    \State 停止原因设为约束冲突，终止循环
  \EndIf
  \State 以 $\frac{z}{\tau}$ 稳定归一化，按固定并列规则构造累计概率至少为 $\rho$ 的最小核
  \State 核内重新归一化，按当前随机数状态抽取词元 $v$
  \State 将 $v$ 追加到响应和前缀，同步缓存与逻辑位置
  \If{$v$ 属于 EOS 集合}
    \State 停止原因设为自然终止，终止循环
  \EndIf
  \State 更新连续文本状态
  \If{命中停止串}
    \State 按协议处理停止串，记录原因并终止循环
  \EndIf
\EndFor
\State \Return 响应词元、连续解码文本及停止原因
\end{lfsalgorithmic}
贪心模式应使用独立的最大值选择分支，不将 $\tau=0$ 送入除法。Top-k、重复惩罚或其他处理器若参与执行，还必须说明硬约束、温度归一化与核选择的执行顺序。
\end{algorithm}

\subsection{重复的跨步反馈}
一段重复输出进入前缀后，模型可能继续给同类续写较大概率；新产生的重复又成为条件，形成跨步反馈。低温度会在重复词元已是最高分时进一步集中概率，重复的原因则不止于此。贪心和束搜索同样可能保留重复路径，随机采样也可能重复。

重复 $n$ 元组屏蔽可以禁止已经出现的片段再次出现，但会误伤需要复述的名称、代码、公式或合法韵律。按已出现词元次数调整分数是一种软约束，也会改变语义词的使用概率。若所有合法候选都被规则排除，算法必须报告冲突或使用事先声明的回退策略。

还应区分模型重复与系统重复。流式分片重复拼接、错误重试或缓存前缀错配也能产生重复文本，调高温度无法解决这类问题。模型分数、解码规则和传输状态对应不同的原因，应分别定位。

\section{条件生成与缓存的边界}

\subsection{提示、任务边界与模板}
仅解码器条件生成通常把系统说明、用户输入与历史消息序列化为一个前缀，再续写响应。消息角色与边界标记属于模型看到的条件，不是界面装饰。把相同自然语言放入不同角色、删除分隔符或改变响应起始标记，可能产生不同的条件分布。

训练时若只监督响应，应保留提示对响应的可见性，同时将提示的损失权重设为零。测试条件生成时通常比较 $p(y\mid c)$，提示本身的概率 $p(c)$ 不进入响应排名；不同提示长度的总序列损失也无法直接作为响应质量比较。

在编码器--解码器模型中，条件 $c$ 先编码为记忆，目标侧仍通过教师强制或逐步生成处理。交叉注意力通常允许读取全部有效源记忆，目标自注意力保持因果。两种结构的分词器、特殊标记与权重布局各不相同，输入输出形状一致并不说明可以互换。

\subsection{缓存复用的分布等价性}
如果权重、位置和可见性不变，新增未来词元不会改变历史位置的隐藏状态。因此历史各层键和值可以保存，下一步只计算新增位置并读取缓存。缓存复用应产生与重新计算相同前缀相同的数学条件分布；它保存的仍是逐位置的键值，不是把原始长前缀压缩为一个固定维数向量。

如果前缀修改、位置方案变化、模型参数更换或掩码不再允许原来的信息范围，缓存便可能失效。滑动窗口也必须明确舍弃哪些历史状态以及如何解释位置。以错误复用换取性能会改变概率模型。

\subsection{生成复现条件与概率语义}
同一权重和提示不足以确定随机采样结果，还需要解码配置、随机数状态、输入模板和执行环境。固定随机种子有利于同一条件下复核，跨设备、后端或数值精度的逐词元一致则另需条件。很小的概率差异也可能使一次离散选择落到不同分支，继而改变整个后续前缀。

词元概率表示模型在给定上下文中的条件偏好，与某个事实命题的置信度是两回事。对多个词元组成的回答，其中一个高概率词元也构不成整段真实性的证据。最终评价应依据任务要求，并保持模型原始分布、解码后分布和实际返回文本的区别。

\begin{note}[生成质量来自三个相互独立的定义]
模型定义 $p_\theta$，解码器定义如何选择或抽样，停止协议定义输出在哪里结束。训练损失验证条件预测，采样规则决定实际输出分布，序列任务评价判断结果是否有用。三者分别描述概率拟合、输出选择和任务效果。
\end{note}


\chapterreferences
\lfsendchapterdocument
